% 高一16_格致中学_周末作业2.tex —— 高一上数学「2026-2027 年格致中学高一上周末作业 2」（试卷版/答案版）
% 试卷版：xelatex 高一16_格致中学_周末作业2.tex
% 答案版：xelatex "\def\isanswer{1}% 高一16_格致中学_周末作业2.tex —— 高一上数学「2026-2027 年格致中学高一上周末作业 2」（试卷版/答案版）
% 试卷版：xelatex 高一16_格致中学_周末作业2.tex
% 答案版：xelatex "\def\isanswer{1}% 高一16_格致中学_周末作业2.tex —— 高一上数学「2026-2027 年格致中学高一上周末作业 2」（试卷版/答案版）
% 试卷版：xelatex 高一16_格致中学_周末作业2.tex
% 答案版：xelatex "\def\isanswer{1}% 高一16_格致中学_周末作业2.tex —— 高一上数学「2026-2027 年格致中学高一上周末作业 2」（试卷版/答案版）
% 试卷版：xelatex 高一16_格致中学_周末作业2.tex
% 答案版：xelatex "\def\isanswer{1}\input{高一16_格致中学_周末作业2.tex}"
%
% 原始材料是微信公众号图文（公众号「破晓」，原文链接 https://mp.weixin.qq.com/s/GcU0DDwonIGn8iVGwVRmpg），
% 正文为 7 张 PNG 页；原文本身就是一份**讲评版**——题目下方直接印出【答案】或【解析】。
% 试题文字、答案与解析均照原卷录入，试卷版把所有行间答案改排 \ans{}、解析收进 \ifanswer…\fi 块。
%
% 与原卷的排版差异（均已在 README「说明」中交代）：
%   ① 原文自**第 3 题**起（第 1、2 题未在该文中发布），故本稿题号亦自 3 起，与高一03（同校作业 1）一致；
%   ② 原卷本就印有「一、填空题」「二、选择题」「三、解答题」三节标题（分见原图 00/02/03.png），
%      本稿照原卷分节，只把标题改排成全稿统一的「大节标题」样式；
%   ③ 原卷把答案直接印在题目下方（选择题第 12–14 题印在【解析】里），本稿统一改为试卷版隐去、
%      答案版以【答案】给出，使两版彻底分离；选择题的括注一律改排「（\quad）」；
%   ④ 原卷的实数集、自然数集符号印作普通斜体 $R$、$N$、$N^*$（第 9、10、12、14、18 题等），
%      本稿按全稿惯例统一写作 $\R$、$\N$、$\N^*$；原卷中文括注用**半角**括号（第 11–14 题的
%      「(   )」），本稿按全稿惯例排全角；
%   ⑤ 第 11 题原卷题干下方本应有一幅三集合 Venn 图，但公众号原图该处为空白（题干与选项之间
%      只有正常行距，逐像素核对过），图确实缺失。本稿照录题干、不补图、不代拟答案，详见该题下的注。
%
% 原卷笔误四处（已逐处放大到能看清字形后核对，本稿按文意订正，并在 README 中写明原委）：
%   ① 第 9 题【解析】③ 原卷印作「③$a=-2$ 时」（与 ② 同为 $a=-2$），但该行紧随的
%      「化为 $1=x-2$，解得 $x=3$」只在 $a=2$ 时成立，且本段结论「$a$ 取值的集合为
%      $\{\pm 2,-\frac{17}{4}\}$」含 $+2$，显系笔误，本稿订正为「③$a=2$ 时」；
%   ② 第 16 题【解析】(1) 原卷印作「若命题……是真命题，所以 $\Delta=m^2-4(2m+5)>0$，
%      解得 $m>10$ 或 $m<-2$」，与题干「是假命题」及紧随的结论「故 $A=\{m\mid -2\leqslant m\leqslant 10\}$」
%      都矛盾（按 $\Delta>0$ 应得 $m>10$ 或 $m<-2$），按文意订正为「是假命题」「$\Delta\leqslant 0$」
%      「$-2\leqslant m\leqslant 10$」（结论不变）。
%   ③ 第 12 题【解析】② 原卷印作「$xy\leqslant 0$ 是 $|x-y|=|x|+|y|$ 的充要条件，因此②是假命题」，
%      但 $xy\leqslant 0$ 与 $|x-y|=|x|+|y|$ 确为充要条件、该命题应为真，与同行的「因此②是假命题」
%      及段末「故选 B」（须恰有 1 个真命题）冲突；而题干 ② 印的正是 $xy<0$（充分不必要条件，恰为假），
%      故本稿按题干与文意把该行订正为「$xy<0$ 是 $|x-y|=|x|+|y|$ 的充分不必要条件」；
%   ④ 第 18 题【解析】(1) 原卷印作「由于 $10=1\times10\times2\times5$」（三个乘号），但该式左端为 10
%      而右端为 100；紧随的两组方程组恰是按 $1\times10$ 与 $2\times5$ 这两个分解写出的，显系漏排等号，
%      本稿订正为「$10=1\times10=2\times5$」。
% 另：第 16 题(2)【解析】的「所以 $A\subset B$」照原卷录入（原卷为**真包含**，与第 10 题①的
%     $A\subseteq I$ 逐处放大比对后确认不同）；2026-09-15 回原图复核时曾误作 $\subseteq$，已订正。
% 原卷存疑两处（无法确证，本稿照抄并加注）：
%   ① 第 8 题只印「【解析】64」；② 第 18(3)【解析】中「反之当 $B=\{x\mid x=2k+1,k\in\Z\}$」。

\documentclass[a4paper,11pt]{article}
\usepackage{common}

\begin{document}

\examtitlest[3]{2026--2027 年格致中学高一上周末作业 2}{集合与常用逻辑用语}

\bigsec{一、填空题}

\begin{enumerate}
\setcounter{enumi}{2}

\item 若非空集合 $A=\{x\mid -2m+6<x<m-2\}$，$B=\{x\mid -m<x<m\}$，且 $A\subset B$，则实数 $m$ 的取值范围是 \blankline[4.4em].

\ans{$\left(\dfrac83,6\right]$}

\item 已知条件 $p:2k-1\leqslant x\leqslant -3k$，条件 $q:-1<x<3$，$p$ 是 $q$ 的必要条件，则实数 $k$ 的取值范围为 \blankline[4.4em].

\ans{$(-\infty,-1]$}

\ifanswer
\solbegin
\step{若 $p$ 是 $q$ 的必要条件，则 $q\Rightarrow p$，即 $\begin{cases}-3k\geqslant 3\\ 2k-1\leqslant -1\end{cases}$，得 $\begin{cases}k\leqslant -1\\ k\leqslant 0\end{cases}$，得 $k\leqslant -1$，}
\step{即实数 $k$ 的取值范围是 $(-\infty,-1]$.}
\solend

\fi

\item 集合 $A=\{0,1,a\}$，$B=(0,2)$，若 $A\cap B=\{1\}$，则实数 $a$ 的取值范围是 \blankline[4.4em].

\ans{$(-\infty,0]\cup[2,+\infty)$}

\item 已知 $\varnothing\subset M\subseteq\N$，且当 $x\in M$ 时，$6-x\in M$，若集合 $M$ 只有三个元素，这样的集合 $M$ 为 \blankline[4.4em].

% 注：原卷第 6 题此处印作「$\varnothing\subset M\subseteq N$」，其中「$N$」为普通斜体 $N$（无空心双竖线）。
%     由答案 $\{0,3,6\}$、$\{1,3,5\}$、$\{2,3,4\}$ 反推，$6-x\in M$ 须把元素限制在 $[0,6]$ 内，
%     故该符号即自然数集 $\N$，本稿按全稿惯例写作 $\N$。
\ans{$\{0,3,6\}$ 或 $\{1,3,5\}$ 或 $\{2,3,4\}$}

\item 任意两个正整数 $x$、$y$，定义某种运算 $\otimes$：$x\otimes y=\begin{cases}x+y & (x\text{ 与 }y\text{ 奇偶相同})\\ x\times y & (x\text{ 与 }y\text{ 奇偶不同})\end{cases}$，则集合 $M=\{(x,y)\mid x\otimes y=6,\ x,y\in\N^*\}$ 中元素的个数是 \blankline[2.4em].

\ans{$9$}

\ifanswer
\solbegin
\step{①当 $x$ 与 $y$ 都为奇数时，有 $1+5=6$，$3+3=6$，}
\step{得 $(1,5)$、$(5,1)$、$(3,3)$，$3$ 个点符合题意，}
\step{②当 $x$ 与 $y$ 都为偶数时，有 $2+4=6$，}
\step{得 $(2,4)$、$(4,2)$，$2$ 个点符合题意，}
\step{③当 $x$ 与 $y$ 一奇一偶时，$1\times 6=6$，$2\times 3=6$，}
\step{得 $(1,6)$、$(6,1)$、$(2,3)$、$(3,2)$，$4$ 个点符合题意，}
\step{所以共有 $9$ 个点符合题意.}
\solend

\fi

\item 已知集合 $A=\{a\mid a=2^x3^y5^z,\ x,y,z\in\N\}$，$B=\{b\mid b\in A,\ 1\leqslant b\leqslant 10\}$，试求集合 $B$ 中含有元素 $1$ 和 $10$ 的真子集的个数为 \blankline[4.4em].

% 注：原卷第 8 题只印「【解析】64」，无过程。按 $B=\{1,2,3,4,5,6,8,9,10\}$（9 个元素）计算，
%     含 $1$ 与 $10$ 的真子集应为 $2^7-1=127$，与 $64$（$=2^6$）不符；原卷口径无法确证，
%     本稿如原卷照录 $64$，不代为改写。
\ans{$64$}

\item 已知 $A=\left\{x\mid \dfrac{x+a}{x^2-4}=1,\ x\in\R\right\}$ 仅有两个子集，则实数 $a$ 取值的集合为 \blankline[4.4em].

\ans{$\left\{\pm 2,-\dfrac{17}{4}\right\}$}

\ifanswer
\solbegin
\step{因为 $A=\{x\mid \dfrac{x+a}{x^2-4}=1,\ x\in\R\}$ 仅有两个子集，所以 $A$ 只有一个元素，$\dfrac{x+a}{x^2-4}=1$，}
\step{①$a\neq\pm 2$，得 $x^2-x-a-4=0$，所以此方程有两个相等的实数根，}
\step{得 $\Delta=1-4(-a-4)=0$，解得 $a=-\dfrac{17}{4}$. 经过验证满足条件.}
\step{②$a=-2$ 时，$\dfrac{x+a}{x^2-4}=1$，化为 $1=x+2$，解得 $x=-1$ 满足条件.}
% 注：原卷此处 ③ 印作「③$a=-2$ 时」（与 ② 同为 $a=-2$），但该行紧随的「化为 $1=x-2$，解得 $x=3$」
%     只在 $a=2$ 时成立（$a=-2$ 时 $(x-2)/(x^2-4)=1/(x+2)=1$ 解得 $x=-1$，即 ② 的情形），
%     且本段结论「$a$ 取值的集合为 $\{\pm 2,-\frac{17}{4}\}$」含 $+2$，显系笔误，按文意订正为 $a=2$.
\step{③$a=2$ 时，$\dfrac{x+a}{x^2-4}=1$，化为 $1=x-2$，解得 $x=3$ 满足条件.}
\step{综上，实数 $a$ 取值的集合为 $\left\{\pm 2,-\dfrac{17}{4}\right\}$.}
\solend

\fi

\item 设集合 $I=\{1,2,3,4,5\}$，若非空集合 $A$ 同时满足①$A\subseteq I$，②$|A|\leqslant \min(A)$（其中 $|A|$ 表示 $A$ 中元素的个数，$\min(A)$ 表示集合 $A$ 中最小元素），称集合 $A$ 为 $I$ 的一个好子集，$I$ 的所有好子集的个数为 \blankline[2.4em].

\ans{$12$}

\ifanswer
\solbegin
\step{①当 $|A|=1$，即集合 $A$ 中元素的个数为 $1$ 时，}
\step{$A$ 的可能情况为 $\{1\}$，$\{2\}$，$\{3\}$，$\{4\}$，$\{5\}$，}
\step{②当 $|A|=2$，即集合 $A$ 中元素的个数为 $2$ 时，}
\step{$A$ 的可能情况为 $\{2,3\}$，$\{2,4\}$，$\{2,5\}$，$\{3,4\}$，$\{3,5\}$，$\{4,5\}$，}
\step{③当 $|A|=3$，即集合 $A$ 中元素的个数为 $3$ 时，$A$ 的可能情况为 $\{3,4,5\}$，}
\step{所以 $I$ 的所有好子集的个数为 $12$.}
\solend

\fi

\end{enumerate}

\bigsec{二、选择题}

\begin{enumerate}
\setcounter{enumi}{10}

% 注：原卷第 11 题的 Venn 图在公众号原图里缺失（题干与选项之间只有正常行距，无任何图形像素），
%     图缺失即无从判断阴影区域，故本稿照录题干与四个选项，不补图、不代拟答案。
\item 下图中的阴影部分的正确表示是（\quad）

A. $(A\cup B)\cap(A\cup C)$\qquad B. $(A\cup B)\cap C$

C. $(A\cup C)\cap(B\cup C)$\qquad D. $(A\cup B)\cap(B\cup C)$

\ans{原卷此处应有三集合 Venn 图，但公众号原图缺失，无从判断阴影区域，故本稿不给出答案}

\item 有下列三个命题：①“$x+y\neq 4$”是“$x\neq 1$ 且 $y\neq 3$”的必要非充分条件；②$xy<0$ 是 $|x-y|=|x|+|y|$ 的充要条件；③已知 $m$、$n\in\Z$，则 $m^2+n^2<5$ 是 $m+n\leqslant 2$ 的充分不必要条件；其中的真命题有（\quad）

A. $0$ 个\qquad B. $1$ 个\qquad C. $2$ 个\qquad D. $3$ 个

\ans{$B$}

\ifanswer
\solbegin
\step{①“$x+y\neq 4$”，推不出“$x\neq 1$ 且 $y\neq 3$”，反之不成立，例如 $x=5$，$y=-1$，则 $x+y=4$，}
\step{因此“$x+y\neq 4$”是“$x\neq 1$ 且 $y\neq 3$”的既不充分也不必要条件，因此①是假命题，}
% 注：原卷 ② 一行印作「②$xy\leqslant 0$ 是 $|x-y|=|x|+|y|$ 的充要条件，因此②是假命题」。
%     但 $xy\leqslant 0$ 与 $|x-y|=|x|+|y|$ 确实互为充要条件，该命题应为**真**，与同行的
%     「因此②是假命题」以及本段末的「故选 B」（须恰有 1 个真命题，即只能 ③ 为真）都冲突；
%     而题干 ② 印的是 $xy<0$，$xy<0$ 只是充分不必要条件，恰为假命题。故本稿按题干与文意
%     把解析该行订正为「$xy<0$ 是 $|x-y|=|x|+|y|$ 的充分不必要条件」。
\step{②$xy<0$ 是 $|x-y|=|x|+|y|$ 的充分不必要条件，因此②是假命题；}
\step{③若 $m$、$n\in\Z$，$m^2+n^2<5$，则 $(m,n)$ 为 $(1,0)$，$(2,0)$，$(1,1)$，$(-1,0)$，$(-2,0)$，$(-1,1)$，$(0,1)$，$(0,2)$，$(0,-1)$，$(0,-2)$，$(-1,-1)$，$(1,-1)$，共 $9$ 组，}
\step{都满足 $m+n\leqslant 2$，反之不成立，例如：$m=3$，$n=-1$，}
\step{所以 $m$、$n\in\Z$，则 $m^2+n^2<5$ 是 $m+n\leqslant 2$ 的充分不必要条件，因此③是真命题.}
\step{故选 $B$.}
\solend

\fi

\item 若 $a$、$b$、$c$ 是 $\triangle ABC$ 的三条边，则“$a^2+b^2+c^2=ab+bc+ca$”是“$\triangle ABC$ 是等腰三角形”的（\quad）

A. 充分不必要条件\qquad B. 必要不充分条件

C. 充要条件\qquad D. 既不充分也不必要条件

\ans{$A$}

\ifanswer
\solbegin
\step{若“$\triangle ABC$ 是等腰三角形”，则当 $a=b\neq c$，}
\step{则 $a^2+b^2+c^2=ab+bc+ca$ 不一定成立，}
\step{若 $a^2+b^2+c^2=ab+bc+ca$，则 $2a^2+2b^2+2c^2=2ab+2bc+2ca$，}
\step{即 $(a-b)^2+(b-c)^2+(c-a)^2=0$，即 $a-b=0$，$b-c=0$，$c-a=0$，}
\step{则 $a=b=c$，则“$\triangle ABC$ 是等腰三角形”成立，}
\step{即“$a^2+b^2+c^2=ab+bc+ca$”是“$\triangle ABC$ 是等腰三角形”充分不必要条件，}
\step{故选 $A$.}
\solend

\fi

\item 设集合 $P_1=\{x\mid x^2+ax+1>0\}$，$P_2=\{x\mid x^2+ax+2>0\}(a\in\R)$，下列说法正确的是（\quad）

A. 对任意 $a$，$P_1$ 是 $P_2$ 的真子集\qquad B. 对任意 $a$，$P_1$ 不是 $P_2$ 的子集

C. 存在 $a$，使得 $P_1$ 不是 $P_2$ 的子集\qquad D. 存在 $a$，使得 $P_2$ 是 $P_1$ 的真子集

\ans{$A$}

\ifanswer
\solbegin
\step{由 $x^2+ax+1>0$ 得 $x^2+ax+2=x^2+ax+1+1>1>0$ 一定成立，}
\step{但当 $x^2+ax+2>0$ 时，$x^2+ax+1>0$ 不一定成立，}
\step{故对于任意的 $a$ 都有 $P_1\bsub P_2$. 故选 $A$.}
\solend

\fi

\end{enumerate}

\bigsec{三、解答题}

\begin{enumerate}
\setcounter{enumi}{14}

\item 求证：关于 $x$ 的方程 $ax^2+bx+c=0$ 有一根为 $1$ 的充要条件是 $a+b+c=0$.

\ans{略}

\ansspace[8]

\item 已知命题“关于 $x$ 的方程 $x^2+mx+2m+5=0$ 有两个不相等的实数根”是假命题.

\begin{enumerate}[label=(\arabic*)]
\item 求实数 $m$ 的取值集合 $A$；
\item 设集合 $B=\{x\mid 1-2a\leqslant x\leqslant a-1\}$，若 $x\in A$ 是 $x\in B$ 的充分不必要条件，求实数 $a$ 的取值范围.
\end{enumerate}

\ans{（1）$A=\{m\mid -2\leqslant m\leqslant 10\}$；（2）$[11,+\infty)$}

\ifanswer
\solbegin
% 注：原卷 (1) 的解析印作「若命题……是真命题，所以 $\Delta=m^2-4(2m+5)>0$，解得 $m>10$ 或 $m<-2$」，
%     与题干「……是假命题」以及紧随的结论「故 $A=\{m\mid -2\leqslant m\leqslant 10\}$」都矛盾
%     （按 $\Delta>0$ 应得 $m>10$ 或 $m<-2$），显系笔误；本稿按文意订正为「是假命题」「$\Delta\leqslant 0$」
%     「$-2\leqslant m\leqslant 10$」，结论不变。
\stepm{(1)}{若命题“关于 $x$ 的方程 $x^2+mx+2m+5=0$ 有两个不相等的实数根”是假命题，所以 $\Delta=m^2-4(2m+5)\leqslant 0$，解得 $-2\leqslant m\leqslant 10$，}
\step{故 $A=\{m\mid -2\leqslant m\leqslant 10\}$.}
\stepm{(2)}{因为 $A=\{m\mid -2\leqslant m\leqslant 10\}$，$x\in A$ 是 $x\in B$ 的充分不必要条件，所以 $A\subset B$.}
\step{即 $\begin{cases}1-2a\leqslant -2\\ a-1\geqslant 10\end{cases}$ 且等号不同时成立，解得 $a\geqslant 11$，}
\step{所以实数 $a$ 的取值范围为 $[11,+\infty)$.}
\solend

\fi

\ansspace[6]

\item 设非空集合 $S=\{x\mid m\leqslant x\leqslant n\}$ 满足：当 $x\in S$ 时，有 $x^2\in S$.

\begin{enumerate}[label=(\arabic*)]
\item 当 $m=1$ 时，求 $S$；
\item 若 $m=-\dfrac12$，求实数 $n$ 的取值范围.
\end{enumerate}

\ans{（1）$S=\{1\}$；（2）$\left[\dfrac14,1\right]$}

\ifanswer
\solbegin
\stepm{(1)}{若 $m=1$ 时，所以 $S=\{x\mid 1\leqslant x\leqslant n\}$，所以 $\forall x$ 有 $1\leqslant x^2\leqslant n^2$，}
\step{所以 $n\leqslant n^2$，所以 $n\leqslant 1$，所以 $S=\{1\}$；}
\stepm{(2)}{若 $m=-\dfrac12$，}
\step{①$-\dfrac12\leqslant n\leqslant 0$ 时，$\forall x\in S$，有 $n^2\leqslant x^2\leqslant \dfrac14$，所以 $x^2\notin S$，舍去，}
\step{②$0<n\leqslant \dfrac12$ 时，$\forall x\in S$，有 $0\leqslant x^2\leqslant \dfrac14$，所以 $\dfrac14\leqslant n\leqslant \dfrac12$，}
\step{③$n>\dfrac12$ 时，$\forall x\in S$，有 $0\leqslant x^2\leqslant n^2$，所以 $n^2\leqslant n$，故 $\dfrac12<n\leqslant 1$.}
\solend

\fi

\ansspace[6]

% 压轴题：其后已无内容，故作答区全用可丢弃胶水（\ansspace*），并在题目前先量一下版面。
\ansneed{7cm}

\item 已知集合 $A=\{x\mid x=m^2-n^2,\ m$、$n\in\Z\}$.

\begin{enumerate}[label=(\arabic*)]
\item 判断 $8$，$9$，$10$ 是否属于集合 $A$；
\item 已知集合 $B=\{x\mid x=2k+1,k\in\Z\}$，证明：“$x\in A$”是“$x\in B$”的必要非充分条件；
\item 求所有满足集合 $A$ 的偶数，并说明理由.
\end{enumerate}

\ans{（1）$8\in A$，$9\in A$，$10\notin A$；（2）略；（3）所有满足集合 $A$ 的偶数为 $4k(k\in\Z)$}

\ifanswer
\solbegin
\stepm{(1)}{由于 $8=3^2-1^2$，满足集合 $A$ 中元素特征 $x=m^2-n^2$，$m$、$n\in\Z$，所以 $8\in A$，}
\step{由于 $9=5^2-4^2$，满足集合 $A$ 中元素特征 $x=m^2-n^2$，$m$、$n\in\Z$，所以 $9\in A$，}
\step{假设 $10=m^2-n^2$，$m$、$n\in\Z$，}
\step{则 $(|m|+|n|)(|m|-|n|)=10$，且 $|m|+|n|>|m|-|n|>0$，}
% 注：原卷此行印作「由于 $10=1\times10\times2\times5$」（三个乘号），右端为 100 而非 10；
%     紧随的两组方程组正是按 $1\times10$ 与 $2\times5$ 两个分解写出的，显系漏排等号，
%     本稿按文意订正为「$10=1\times10=2\times5$」（详见文件头「原卷笔误」④）。
\step{由于 $10=1\times 10=2\times 5$，所以 $\begin{cases}|m|+|n|=10\\ |m|-|n|=1\end{cases}$ 或 $\begin{cases}|m|+|n|=5\\ |m|-|n|=2\end{cases}$，}
\step{显然均无整数解，所以 $10\notin A$；}
\stepm{(2)}{集合 $B=\{x\mid x=2k+1,k\in\Z\}$，则恒有 $2k+1=(k+1)^2-k^2$，}
\step{即满足集合 $A$ 中元素特征 $x=m^2-n^2$，$m$、$n\in\Z$，所以 $2k+1\in A$，}
\step{即一切奇数都属于 $A$；}
\step{反之满足这个集合中的元素不一定全是奇数，如 $8\in A$，}
\step{所以 $x\in A$ 是 $x\in B$ 的必要非充分条件，}
\stepm{(3)}{集合 $A=\{x\mid x=m^2-n^2\}(m,n\in\Z)$，而 $m^2-n^2=(m+n)(m-n)$，}
\step{①当 $m$ 和 $n$ 一奇一偶时，$m+n$ 和 $m-n$ 均为奇数，}
\step{所以 $(m+n)(m-n)$ 为奇数，不满足题意；}
\step{②当 $m$ 和 $n$ 同为奇数和偶数时，$m-n$，$m+n$ 均为偶数，}
\step{所以 $(m+n)(m-n)$ 为 $4$ 的倍数，}
% 注：原卷此处印作「反之当 $B=\{x\mid x=2k+1,k\in\Z\}$」（与第（2）问定义的 $B$ 一字不差），
%     但紧随的推导「不妨令 $m+n=2k$，$m-n=2$」所得正是 $x=(m+n)(m-n)=4k$，
%     且本段结论为「所有满足集合 $A$ 的偶数为 $4k(k\in\Z)$」；此处 $B$ 疑为 $\{x\mid x=4k,k\in\Z\}$
%     之误（或漏改自第（2）问），但无法确证，本稿照抄原卷写法。
\step{反之当 $B=\{x\mid x=2k+1,k\in\Z\}$，则不妨令 $\begin{cases}m+n=2k\\ m-n=2\end{cases}$，}
\step{可解得 $\begin{cases}m=k+1\\ n=k-1\end{cases}$，满足集合 $A$ 中元素特征 $x=m^2-n^2$，$m$、$n\in\Z$，}
\step{所以满足集合 $A$ 的偶数为 $4k(k\in\Z)$；}
\step{综上所述，所有满足集合 $A$ 的偶数为 $4k(k\in\Z)$.}
\solend

\fi

\ansspace*[10]

\end{enumerate}

\end{document}
"
%
% 原始材料是微信公众号图文（公众号「破晓」，原文链接 https://mp.weixin.qq.com/s/GcU0DDwonIGn8iVGwVRmpg），
% 正文为 7 张 PNG 页；原文本身就是一份**讲评版**——题目下方直接印出【答案】或【解析】。
% 试题文字、答案与解析均照原卷录入，试卷版把所有行间答案改排 \ans{}、解析收进 \ifanswer…\fi 块。
%
% 与原卷的排版差异（均已在 README「说明」中交代）：
%   ① 原文自**第 3 题**起（第 1、2 题未在该文中发布），故本稿题号亦自 3 起，与高一03（同校作业 1）一致；
%   ② 原卷本就印有「一、填空题」「二、选择题」「三、解答题」三节标题（分见原图 00/02/03.png），
%      本稿照原卷分节，只把标题改排成全稿统一的「大节标题」样式；
%   ③ 原卷把答案直接印在题目下方（选择题第 12–14 题印在【解析】里），本稿统一改为试卷版隐去、
%      答案版以【答案】给出，使两版彻底分离；选择题的括注一律改排「（\quad）」；
%   ④ 原卷的实数集、自然数集符号印作普通斜体 $R$、$N$、$N^*$（第 9、10、12、14、18 题等），
%      本稿按全稿惯例统一写作 $\R$、$\N$、$\N^*$；原卷中文括注用**半角**括号（第 11–14 题的
%      「(   )」），本稿按全稿惯例排全角；
%   ⑤ 第 11 题原卷题干下方本应有一幅三集合 Venn 图，但公众号原图该处为空白（题干与选项之间
%      只有正常行距，逐像素核对过），图确实缺失。本稿照录题干、不补图、不代拟答案，详见该题下的注。
%
% 原卷笔误四处（已逐处放大到能看清字形后核对，本稿按文意订正，并在 README 中写明原委）：
%   ① 第 9 题【解析】③ 原卷印作「③$a=-2$ 时」（与 ② 同为 $a=-2$），但该行紧随的
%      「化为 $1=x-2$，解得 $x=3$」只在 $a=2$ 时成立，且本段结论「$a$ 取值的集合为
%      $\{\pm 2,-\frac{17}{4}\}$」含 $+2$，显系笔误，本稿订正为「③$a=2$ 时」；
%   ② 第 16 题【解析】(1) 原卷印作「若命题……是真命题，所以 $\Delta=m^2-4(2m+5)>0$，
%      解得 $m>10$ 或 $m<-2$」，与题干「是假命题」及紧随的结论「故 $A=\{m\mid -2\leqslant m\leqslant 10\}$」
%      都矛盾（按 $\Delta>0$ 应得 $m>10$ 或 $m<-2$），按文意订正为「是假命题」「$\Delta\leqslant 0$」
%      「$-2\leqslant m\leqslant 10$」（结论不变）。
%   ③ 第 12 题【解析】② 原卷印作「$xy\leqslant 0$ 是 $|x-y|=|x|+|y|$ 的充要条件，因此②是假命题」，
%      但 $xy\leqslant 0$ 与 $|x-y|=|x|+|y|$ 确为充要条件、该命题应为真，与同行的「因此②是假命题」
%      及段末「故选 B」（须恰有 1 个真命题）冲突；而题干 ② 印的正是 $xy<0$（充分不必要条件，恰为假），
%      故本稿按题干与文意把该行订正为「$xy<0$ 是 $|x-y|=|x|+|y|$ 的充分不必要条件」；
%   ④ 第 18 题【解析】(1) 原卷印作「由于 $10=1\times10\times2\times5$」（三个乘号），但该式左端为 10
%      而右端为 100；紧随的两组方程组恰是按 $1\times10$ 与 $2\times5$ 这两个分解写出的，显系漏排等号，
%      本稿订正为「$10=1\times10=2\times5$」。
% 另：第 16 题(2)【解析】的「所以 $A\subset B$」照原卷录入（原卷为**真包含**，与第 10 题①的
%     $A\subseteq I$ 逐处放大比对后确认不同）；2026-09-15 回原图复核时曾误作 $\subseteq$，已订正。
% 原卷存疑两处（无法确证，本稿照抄并加注）：
%   ① 第 8 题只印「【解析】64」；② 第 18(3)【解析】中「反之当 $B=\{x\mid x=2k+1,k\in\Z\}$」。

\documentclass[a4paper,11pt]{article}
\usepackage{common}

\begin{document}

\examtitlest[3]{2026--2027 年格致中学高一上周末作业 2}{集合与常用逻辑用语}

\bigsec{一、填空题}

\begin{enumerate}
\setcounter{enumi}{2}

\item 若非空集合 $A=\{x\mid -2m+6<x<m-2\}$，$B=\{x\mid -m<x<m\}$，且 $A\subset B$，则实数 $m$ 的取值范围是 \blankline[4.4em].

\ans{$\left(\dfrac83,6\right]$}

\item 已知条件 $p:2k-1\leqslant x\leqslant -3k$，条件 $q:-1<x<3$，$p$ 是 $q$ 的必要条件，则实数 $k$ 的取值范围为 \blankline[4.4em].

\ans{$(-\infty,-1]$}

\ifanswer
\solbegin
\step{若 $p$ 是 $q$ 的必要条件，则 $q\Rightarrow p$，即 $\begin{cases}-3k\geqslant 3\\ 2k-1\leqslant -1\end{cases}$，得 $\begin{cases}k\leqslant -1\\ k\leqslant 0\end{cases}$，得 $k\leqslant -1$，}
\step{即实数 $k$ 的取值范围是 $(-\infty,-1]$.}
\solend

\fi

\item 集合 $A=\{0,1,a\}$，$B=(0,2)$，若 $A\cap B=\{1\}$，则实数 $a$ 的取值范围是 \blankline[4.4em].

\ans{$(-\infty,0]\cup[2,+\infty)$}

\item 已知 $\varnothing\subset M\subseteq\N$，且当 $x\in M$ 时，$6-x\in M$，若集合 $M$ 只有三个元素，这样的集合 $M$ 为 \blankline[4.4em].

% 注：原卷第 6 题此处印作「$\varnothing\subset M\subseteq N$」，其中「$N$」为普通斜体 $N$（无空心双竖线）。
%     由答案 $\{0,3,6\}$、$\{1,3,5\}$、$\{2,3,4\}$ 反推，$6-x\in M$ 须把元素限制在 $[0,6]$ 内，
%     故该符号即自然数集 $\N$，本稿按全稿惯例写作 $\N$。
\ans{$\{0,3,6\}$ 或 $\{1,3,5\}$ 或 $\{2,3,4\}$}

\item 任意两个正整数 $x$、$y$，定义某种运算 $\otimes$：$x\otimes y=\begin{cases}x+y & (x\text{ 与 }y\text{ 奇偶相同})\\ x\times y & (x\text{ 与 }y\text{ 奇偶不同})\end{cases}$，则集合 $M=\{(x,y)\mid x\otimes y=6,\ x,y\in\N^*\}$ 中元素的个数是 \blankline[2.4em].

\ans{$9$}

\ifanswer
\solbegin
\step{①当 $x$ 与 $y$ 都为奇数时，有 $1+5=6$，$3+3=6$，}
\step{得 $(1,5)$、$(5,1)$、$(3,3)$，$3$ 个点符合题意，}
\step{②当 $x$ 与 $y$ 都为偶数时，有 $2+4=6$，}
\step{得 $(2,4)$、$(4,2)$，$2$ 个点符合题意，}
\step{③当 $x$ 与 $y$ 一奇一偶时，$1\times 6=6$，$2\times 3=6$，}
\step{得 $(1,6)$、$(6,1)$、$(2,3)$、$(3,2)$，$4$ 个点符合题意，}
\step{所以共有 $9$ 个点符合题意.}
\solend

\fi

\item 已知集合 $A=\{a\mid a=2^x3^y5^z,\ x,y,z\in\N\}$，$B=\{b\mid b\in A,\ 1\leqslant b\leqslant 10\}$，试求集合 $B$ 中含有元素 $1$ 和 $10$ 的真子集的个数为 \blankline[4.4em].

% 注：原卷第 8 题只印「【解析】64」，无过程。按 $B=\{1,2,3,4,5,6,8,9,10\}$（9 个元素）计算，
%     含 $1$ 与 $10$ 的真子集应为 $2^7-1=127$，与 $64$（$=2^6$）不符；原卷口径无法确证，
%     本稿如原卷照录 $64$，不代为改写。
\ans{$64$}

\item 已知 $A=\left\{x\mid \dfrac{x+a}{x^2-4}=1,\ x\in\R\right\}$ 仅有两个子集，则实数 $a$ 取值的集合为 \blankline[4.4em].

\ans{$\left\{\pm 2,-\dfrac{17}{4}\right\}$}

\ifanswer
\solbegin
\step{因为 $A=\{x\mid \dfrac{x+a}{x^2-4}=1,\ x\in\R\}$ 仅有两个子集，所以 $A$ 只有一个元素，$\dfrac{x+a}{x^2-4}=1$，}
\step{①$a\neq\pm 2$，得 $x^2-x-a-4=0$，所以此方程有两个相等的实数根，}
\step{得 $\Delta=1-4(-a-4)=0$，解得 $a=-\dfrac{17}{4}$. 经过验证满足条件.}
\step{②$a=-2$ 时，$\dfrac{x+a}{x^2-4}=1$，化为 $1=x+2$，解得 $x=-1$ 满足条件.}
% 注：原卷此处 ③ 印作「③$a=-2$ 时」（与 ② 同为 $a=-2$），但该行紧随的「化为 $1=x-2$，解得 $x=3$」
%     只在 $a=2$ 时成立（$a=-2$ 时 $(x-2)/(x^2-4)=1/(x+2)=1$ 解得 $x=-1$，即 ② 的情形），
%     且本段结论「$a$ 取值的集合为 $\{\pm 2,-\frac{17}{4}\}$」含 $+2$，显系笔误，按文意订正为 $a=2$.
\step{③$a=2$ 时，$\dfrac{x+a}{x^2-4}=1$，化为 $1=x-2$，解得 $x=3$ 满足条件.}
\step{综上，实数 $a$ 取值的集合为 $\left\{\pm 2,-\dfrac{17}{4}\right\}$.}
\solend

\fi

\item 设集合 $I=\{1,2,3,4,5\}$，若非空集合 $A$ 同时满足①$A\subseteq I$，②$|A|\leqslant \min(A)$（其中 $|A|$ 表示 $A$ 中元素的个数，$\min(A)$ 表示集合 $A$ 中最小元素），称集合 $A$ 为 $I$ 的一个好子集，$I$ 的所有好子集的个数为 \blankline[2.4em].

\ans{$12$}

\ifanswer
\solbegin
\step{①当 $|A|=1$，即集合 $A$ 中元素的个数为 $1$ 时，}
\step{$A$ 的可能情况为 $\{1\}$，$\{2\}$，$\{3\}$，$\{4\}$，$\{5\}$，}
\step{②当 $|A|=2$，即集合 $A$ 中元素的个数为 $2$ 时，}
\step{$A$ 的可能情况为 $\{2,3\}$，$\{2,4\}$，$\{2,5\}$，$\{3,4\}$，$\{3,5\}$，$\{4,5\}$，}
\step{③当 $|A|=3$，即集合 $A$ 中元素的个数为 $3$ 时，$A$ 的可能情况为 $\{3,4,5\}$，}
\step{所以 $I$ 的所有好子集的个数为 $12$.}
\solend

\fi

\end{enumerate}

\bigsec{二、选择题}

\begin{enumerate}
\setcounter{enumi}{10}

% 注：原卷第 11 题的 Venn 图在公众号原图里缺失（题干与选项之间只有正常行距，无任何图形像素），
%     图缺失即无从判断阴影区域，故本稿照录题干与四个选项，不补图、不代拟答案。
\item 下图中的阴影部分的正确表示是（\quad）

A. $(A\cup B)\cap(A\cup C)$\qquad B. $(A\cup B)\cap C$

C. $(A\cup C)\cap(B\cup C)$\qquad D. $(A\cup B)\cap(B\cup C)$

\ans{原卷此处应有三集合 Venn 图，但公众号原图缺失，无从判断阴影区域，故本稿不给出答案}

\item 有下列三个命题：①“$x+y\neq 4$”是“$x\neq 1$ 且 $y\neq 3$”的必要非充分条件；②$xy<0$ 是 $|x-y|=|x|+|y|$ 的充要条件；③已知 $m$、$n\in\Z$，则 $m^2+n^2<5$ 是 $m+n\leqslant 2$ 的充分不必要条件；其中的真命题有（\quad）

A. $0$ 个\qquad B. $1$ 个\qquad C. $2$ 个\qquad D. $3$ 个

\ans{$B$}

\ifanswer
\solbegin
\step{①“$x+y\neq 4$”，推不出“$x\neq 1$ 且 $y\neq 3$”，反之不成立，例如 $x=5$，$y=-1$，则 $x+y=4$，}
\step{因此“$x+y\neq 4$”是“$x\neq 1$ 且 $y\neq 3$”的既不充分也不必要条件，因此①是假命题，}
% 注：原卷 ② 一行印作「②$xy\leqslant 0$ 是 $|x-y|=|x|+|y|$ 的充要条件，因此②是假命题」。
%     但 $xy\leqslant 0$ 与 $|x-y|=|x|+|y|$ 确实互为充要条件，该命题应为**真**，与同行的
%     「因此②是假命题」以及本段末的「故选 B」（须恰有 1 个真命题，即只能 ③ 为真）都冲突；
%     而题干 ② 印的是 $xy<0$，$xy<0$ 只是充分不必要条件，恰为假命题。故本稿按题干与文意
%     把解析该行订正为「$xy<0$ 是 $|x-y|=|x|+|y|$ 的充分不必要条件」。
\step{②$xy<0$ 是 $|x-y|=|x|+|y|$ 的充分不必要条件，因此②是假命题；}
\step{③若 $m$、$n\in\Z$，$m^2+n^2<5$，则 $(m,n)$ 为 $(1,0)$，$(2,0)$，$(1,1)$，$(-1,0)$，$(-2,0)$，$(-1,1)$，$(0,1)$，$(0,2)$，$(0,-1)$，$(0,-2)$，$(-1,-1)$，$(1,-1)$，共 $9$ 组，}
\step{都满足 $m+n\leqslant 2$，反之不成立，例如：$m=3$，$n=-1$，}
\step{所以 $m$、$n\in\Z$，则 $m^2+n^2<5$ 是 $m+n\leqslant 2$ 的充分不必要条件，因此③是真命题.}
\step{故选 $B$.}
\solend

\fi

\item 若 $a$、$b$、$c$ 是 $\triangle ABC$ 的三条边，则“$a^2+b^2+c^2=ab+bc+ca$”是“$\triangle ABC$ 是等腰三角形”的（\quad）

A. 充分不必要条件\qquad B. 必要不充分条件

C. 充要条件\qquad D. 既不充分也不必要条件

\ans{$A$}

\ifanswer
\solbegin
\step{若“$\triangle ABC$ 是等腰三角形”，则当 $a=b\neq c$，}
\step{则 $a^2+b^2+c^2=ab+bc+ca$ 不一定成立，}
\step{若 $a^2+b^2+c^2=ab+bc+ca$，则 $2a^2+2b^2+2c^2=2ab+2bc+2ca$，}
\step{即 $(a-b)^2+(b-c)^2+(c-a)^2=0$，即 $a-b=0$，$b-c=0$，$c-a=0$，}
\step{则 $a=b=c$，则“$\triangle ABC$ 是等腰三角形”成立，}
\step{即“$a^2+b^2+c^2=ab+bc+ca$”是“$\triangle ABC$ 是等腰三角形”充分不必要条件，}
\step{故选 $A$.}
\solend

\fi

\item 设集合 $P_1=\{x\mid x^2+ax+1>0\}$，$P_2=\{x\mid x^2+ax+2>0\}(a\in\R)$，下列说法正确的是（\quad）

A. 对任意 $a$，$P_1$ 是 $P_2$ 的真子集\qquad B. 对任意 $a$，$P_1$ 不是 $P_2$ 的子集

C. 存在 $a$，使得 $P_1$ 不是 $P_2$ 的子集\qquad D. 存在 $a$，使得 $P_2$ 是 $P_1$ 的真子集

\ans{$A$}

\ifanswer
\solbegin
\step{由 $x^2+ax+1>0$ 得 $x^2+ax+2=x^2+ax+1+1>1>0$ 一定成立，}
\step{但当 $x^2+ax+2>0$ 时，$x^2+ax+1>0$ 不一定成立，}
\step{故对于任意的 $a$ 都有 $P_1\bsub P_2$. 故选 $A$.}
\solend

\fi

\end{enumerate}

\bigsec{三、解答题}

\begin{enumerate}
\setcounter{enumi}{14}

\item 求证：关于 $x$ 的方程 $ax^2+bx+c=0$ 有一根为 $1$ 的充要条件是 $a+b+c=0$.

\ans{略}

\ansspace[8]

\item 已知命题“关于 $x$ 的方程 $x^2+mx+2m+5=0$ 有两个不相等的实数根”是假命题.

\begin{enumerate}[label=(\arabic*)]
\item 求实数 $m$ 的取值集合 $A$；
\item 设集合 $B=\{x\mid 1-2a\leqslant x\leqslant a-1\}$，若 $x\in A$ 是 $x\in B$ 的充分不必要条件，求实数 $a$ 的取值范围.
\end{enumerate}

\ans{（1）$A=\{m\mid -2\leqslant m\leqslant 10\}$；（2）$[11,+\infty)$}

\ifanswer
\solbegin
% 注：原卷 (1) 的解析印作「若命题……是真命题，所以 $\Delta=m^2-4(2m+5)>0$，解得 $m>10$ 或 $m<-2$」，
%     与题干「……是假命题」以及紧随的结论「故 $A=\{m\mid -2\leqslant m\leqslant 10\}$」都矛盾
%     （按 $\Delta>0$ 应得 $m>10$ 或 $m<-2$），显系笔误；本稿按文意订正为「是假命题」「$\Delta\leqslant 0$」
%     「$-2\leqslant m\leqslant 10$」，结论不变。
\stepm{(1)}{若命题“关于 $x$ 的方程 $x^2+mx+2m+5=0$ 有两个不相等的实数根”是假命题，所以 $\Delta=m^2-4(2m+5)\leqslant 0$，解得 $-2\leqslant m\leqslant 10$，}
\step{故 $A=\{m\mid -2\leqslant m\leqslant 10\}$.}
\stepm{(2)}{因为 $A=\{m\mid -2\leqslant m\leqslant 10\}$，$x\in A$ 是 $x\in B$ 的充分不必要条件，所以 $A\subset B$.}
\step{即 $\begin{cases}1-2a\leqslant -2\\ a-1\geqslant 10\end{cases}$ 且等号不同时成立，解得 $a\geqslant 11$，}
\step{所以实数 $a$ 的取值范围为 $[11,+\infty)$.}
\solend

\fi

\ansspace[6]

\item 设非空集合 $S=\{x\mid m\leqslant x\leqslant n\}$ 满足：当 $x\in S$ 时，有 $x^2\in S$.

\begin{enumerate}[label=(\arabic*)]
\item 当 $m=1$ 时，求 $S$；
\item 若 $m=-\dfrac12$，求实数 $n$ 的取值范围.
\end{enumerate}

\ans{（1）$S=\{1\}$；（2）$\left[\dfrac14,1\right]$}

\ifanswer
\solbegin
\stepm{(1)}{若 $m=1$ 时，所以 $S=\{x\mid 1\leqslant x\leqslant n\}$，所以 $\forall x$ 有 $1\leqslant x^2\leqslant n^2$，}
\step{所以 $n\leqslant n^2$，所以 $n\leqslant 1$，所以 $S=\{1\}$；}
\stepm{(2)}{若 $m=-\dfrac12$，}
\step{①$-\dfrac12\leqslant n\leqslant 0$ 时，$\forall x\in S$，有 $n^2\leqslant x^2\leqslant \dfrac14$，所以 $x^2\notin S$，舍去，}
\step{②$0<n\leqslant \dfrac12$ 时，$\forall x\in S$，有 $0\leqslant x^2\leqslant \dfrac14$，所以 $\dfrac14\leqslant n\leqslant \dfrac12$，}
\step{③$n>\dfrac12$ 时，$\forall x\in S$，有 $0\leqslant x^2\leqslant n^2$，所以 $n^2\leqslant n$，故 $\dfrac12<n\leqslant 1$.}
\solend

\fi

\ansspace[6]

% 压轴题：其后已无内容，故作答区全用可丢弃胶水（\ansspace*），并在题目前先量一下版面。
\ansneed{7cm}

\item 已知集合 $A=\{x\mid x=m^2-n^2,\ m$、$n\in\Z\}$.

\begin{enumerate}[label=(\arabic*)]
\item 判断 $8$，$9$，$10$ 是否属于集合 $A$；
\item 已知集合 $B=\{x\mid x=2k+1,k\in\Z\}$，证明：“$x\in A$”是“$x\in B$”的必要非充分条件；
\item 求所有满足集合 $A$ 的偶数，并说明理由.
\end{enumerate}

\ans{（1）$8\in A$，$9\in A$，$10\notin A$；（2）略；（3）所有满足集合 $A$ 的偶数为 $4k(k\in\Z)$}

\ifanswer
\solbegin
\stepm{(1)}{由于 $8=3^2-1^2$，满足集合 $A$ 中元素特征 $x=m^2-n^2$，$m$、$n\in\Z$，所以 $8\in A$，}
\step{由于 $9=5^2-4^2$，满足集合 $A$ 中元素特征 $x=m^2-n^2$，$m$、$n\in\Z$，所以 $9\in A$，}
\step{假设 $10=m^2-n^2$，$m$、$n\in\Z$，}
\step{则 $(|m|+|n|)(|m|-|n|)=10$，且 $|m|+|n|>|m|-|n|>0$，}
% 注：原卷此行印作「由于 $10=1\times10\times2\times5$」（三个乘号），右端为 100 而非 10；
%     紧随的两组方程组正是按 $1\times10$ 与 $2\times5$ 两个分解写出的，显系漏排等号，
%     本稿按文意订正为「$10=1\times10=2\times5$」（详见文件头「原卷笔误」④）。
\step{由于 $10=1\times 10=2\times 5$，所以 $\begin{cases}|m|+|n|=10\\ |m|-|n|=1\end{cases}$ 或 $\begin{cases}|m|+|n|=5\\ |m|-|n|=2\end{cases}$，}
\step{显然均无整数解，所以 $10\notin A$；}
\stepm{(2)}{集合 $B=\{x\mid x=2k+1,k\in\Z\}$，则恒有 $2k+1=(k+1)^2-k^2$，}
\step{即满足集合 $A$ 中元素特征 $x=m^2-n^2$，$m$、$n\in\Z$，所以 $2k+1\in A$，}
\step{即一切奇数都属于 $A$；}
\step{反之满足这个集合中的元素不一定全是奇数，如 $8\in A$，}
\step{所以 $x\in A$ 是 $x\in B$ 的必要非充分条件，}
\stepm{(3)}{集合 $A=\{x\mid x=m^2-n^2\}(m,n\in\Z)$，而 $m^2-n^2=(m+n)(m-n)$，}
\step{①当 $m$ 和 $n$ 一奇一偶时，$m+n$ 和 $m-n$ 均为奇数，}
\step{所以 $(m+n)(m-n)$ 为奇数，不满足题意；}
\step{②当 $m$ 和 $n$ 同为奇数和偶数时，$m-n$，$m+n$ 均为偶数，}
\step{所以 $(m+n)(m-n)$ 为 $4$ 的倍数，}
% 注：原卷此处印作「反之当 $B=\{x\mid x=2k+1,k\in\Z\}$」（与第（2）问定义的 $B$ 一字不差），
%     但紧随的推导「不妨令 $m+n=2k$，$m-n=2$」所得正是 $x=(m+n)(m-n)=4k$，
%     且本段结论为「所有满足集合 $A$ 的偶数为 $4k(k\in\Z)$」；此处 $B$ 疑为 $\{x\mid x=4k,k\in\Z\}$
%     之误（或漏改自第（2）问），但无法确证，本稿照抄原卷写法。
\step{反之当 $B=\{x\mid x=2k+1,k\in\Z\}$，则不妨令 $\begin{cases}m+n=2k\\ m-n=2\end{cases}$，}
\step{可解得 $\begin{cases}m=k+1\\ n=k-1\end{cases}$，满足集合 $A$ 中元素特征 $x=m^2-n^2$，$m$、$n\in\Z$，}
\step{所以满足集合 $A$ 的偶数为 $4k(k\in\Z)$；}
\step{综上所述，所有满足集合 $A$ 的偶数为 $4k(k\in\Z)$.}
\solend

\fi

\ansspace*[10]

\end{enumerate}

\end{document}
"
%
% 原始材料是微信公众号图文（公众号「破晓」，原文链接 https://mp.weixin.qq.com/s/GcU0DDwonIGn8iVGwVRmpg），
% 正文为 7 张 PNG 页；原文本身就是一份**讲评版**——题目下方直接印出【答案】或【解析】。
% 试题文字、答案与解析均照原卷录入，试卷版把所有行间答案改排 \ans{}、解析收进 \ifanswer…\fi 块。
%
% 与原卷的排版差异（均已在 README「说明」中交代）：
%   ① 原文自**第 3 题**起（第 1、2 题未在该文中发布），故本稿题号亦自 3 起，与高一03（同校作业 1）一致；
%   ② 原卷本就印有「一、填空题」「二、选择题」「三、解答题」三节标题（分见原图 00/02/03.png），
%      本稿照原卷分节，只把标题改排成全稿统一的「大节标题」样式；
%   ③ 原卷把答案直接印在题目下方（选择题第 12–14 题印在【解析】里），本稿统一改为试卷版隐去、
%      答案版以【答案】给出，使两版彻底分离；选择题的括注一律改排「（\quad）」；
%   ④ 原卷的实数集、自然数集符号印作普通斜体 $R$、$N$、$N^*$（第 9、10、12、14、18 题等），
%      本稿按全稿惯例统一写作 $\R$、$\N$、$\N^*$；原卷中文括注用**半角**括号（第 11–14 题的
%      「(   )」），本稿按全稿惯例排全角；
%   ⑤ 第 11 题原卷题干下方本应有一幅三集合 Venn 图，但公众号原图该处为空白（题干与选项之间
%      只有正常行距，逐像素核对过），图确实缺失。本稿照录题干、不补图、不代拟答案，详见该题下的注。
%
% 原卷笔误四处（已逐处放大到能看清字形后核对，本稿按文意订正，并在 README 中写明原委）：
%   ① 第 9 题【解析】③ 原卷印作「③$a=-2$ 时」（与 ② 同为 $a=-2$），但该行紧随的
%      「化为 $1=x-2$，解得 $x=3$」只在 $a=2$ 时成立，且本段结论「$a$ 取值的集合为
%      $\{\pm 2,-\frac{17}{4}\}$」含 $+2$，显系笔误，本稿订正为「③$a=2$ 时」；
%   ② 第 16 题【解析】(1) 原卷印作「若命题……是真命题，所以 $\Delta=m^2-4(2m+5)>0$，
%      解得 $m>10$ 或 $m<-2$」，与题干「是假命题」及紧随的结论「故 $A=\{m\mid -2\leqslant m\leqslant 10\}$」
%      都矛盾（按 $\Delta>0$ 应得 $m>10$ 或 $m<-2$），按文意订正为「是假命题」「$\Delta\leqslant 0$」
%      「$-2\leqslant m\leqslant 10$」（结论不变）。
%   ③ 第 12 题【解析】② 原卷印作「$xy\leqslant 0$ 是 $|x-y|=|x|+|y|$ 的充要条件，因此②是假命题」，
%      但 $xy\leqslant 0$ 与 $|x-y|=|x|+|y|$ 确为充要条件、该命题应为真，与同行的「因此②是假命题」
%      及段末「故选 B」（须恰有 1 个真命题）冲突；而题干 ② 印的正是 $xy<0$（充分不必要条件，恰为假），
%      故本稿按题干与文意把该行订正为「$xy<0$ 是 $|x-y|=|x|+|y|$ 的充分不必要条件」；
%   ④ 第 18 题【解析】(1) 原卷印作「由于 $10=1\times10\times2\times5$」（三个乘号），但该式左端为 10
%      而右端为 100；紧随的两组方程组恰是按 $1\times10$ 与 $2\times5$ 这两个分解写出的，显系漏排等号，
%      本稿订正为「$10=1\times10=2\times5$」。
% 另：第 16 题(2)【解析】的「所以 $A\subset B$」照原卷录入（原卷为**真包含**，与第 10 题①的
%     $A\subseteq I$ 逐处放大比对后确认不同）；2026-09-15 回原图复核时曾误作 $\subseteq$，已订正。
% 原卷存疑两处（无法确证，本稿照抄并加注）：
%   ① 第 8 题只印「【解析】64」；② 第 18(3)【解析】中「反之当 $B=\{x\mid x=2k+1,k\in\Z\}$」。

\documentclass[a4paper,11pt]{article}
\usepackage{common}

\begin{document}

\examtitlest[3]{2026--2027 年格致中学高一上周末作业 2}{集合与常用逻辑用语}

\bigsec{一、填空题}

\begin{enumerate}
\setcounter{enumi}{2}

\item 若非空集合 $A=\{x\mid -2m+6<x<m-2\}$，$B=\{x\mid -m<x<m\}$，且 $A\subset B$，则实数 $m$ 的取值范围是 \blankline[4.4em].

\ans{$\left(\dfrac83,6\right]$}

\item 已知条件 $p:2k-1\leqslant x\leqslant -3k$，条件 $q:-1<x<3$，$p$ 是 $q$ 的必要条件，则实数 $k$ 的取值范围为 \blankline[4.4em].

\ans{$(-\infty,-1]$}

\ifanswer
\solbegin
\step{若 $p$ 是 $q$ 的必要条件，则 $q\Rightarrow p$，即 $\begin{cases}-3k\geqslant 3\\ 2k-1\leqslant -1\end{cases}$，得 $\begin{cases}k\leqslant -1\\ k\leqslant 0\end{cases}$，得 $k\leqslant -1$，}
\step{即实数 $k$ 的取值范围是 $(-\infty,-1]$.}
\solend

\fi

\item 集合 $A=\{0,1,a\}$，$B=(0,2)$，若 $A\cap B=\{1\}$，则实数 $a$ 的取值范围是 \blankline[4.4em].

\ans{$(-\infty,0]\cup[2,+\infty)$}

\item 已知 $\varnothing\subset M\subseteq\N$，且当 $x\in M$ 时，$6-x\in M$，若集合 $M$ 只有三个元素，这样的集合 $M$ 为 \blankline[4.4em].

% 注：原卷第 6 题此处印作「$\varnothing\subset M\subseteq N$」，其中「$N$」为普通斜体 $N$（无空心双竖线）。
%     由答案 $\{0,3,6\}$、$\{1,3,5\}$、$\{2,3,4\}$ 反推，$6-x\in M$ 须把元素限制在 $[0,6]$ 内，
%     故该符号即自然数集 $\N$，本稿按全稿惯例写作 $\N$。
\ans{$\{0,3,6\}$ 或 $\{1,3,5\}$ 或 $\{2,3,4\}$}

\item 任意两个正整数 $x$、$y$，定义某种运算 $\otimes$：$x\otimes y=\begin{cases}x+y & (x\text{ 与 }y\text{ 奇偶相同})\\ x\times y & (x\text{ 与 }y\text{ 奇偶不同})\end{cases}$，则集合 $M=\{(x,y)\mid x\otimes y=6,\ x,y\in\N^*\}$ 中元素的个数是 \blankline[2.4em].

\ans{$9$}

\ifanswer
\solbegin
\step{①当 $x$ 与 $y$ 都为奇数时，有 $1+5=6$，$3+3=6$，}
\step{得 $(1,5)$、$(5,1)$、$(3,3)$，$3$ 个点符合题意，}
\step{②当 $x$ 与 $y$ 都为偶数时，有 $2+4=6$，}
\step{得 $(2,4)$、$(4,2)$，$2$ 个点符合题意，}
\step{③当 $x$ 与 $y$ 一奇一偶时，$1\times 6=6$，$2\times 3=6$，}
\step{得 $(1,6)$、$(6,1)$、$(2,3)$、$(3,2)$，$4$ 个点符合题意，}
\step{所以共有 $9$ 个点符合题意.}
\solend

\fi

\item 已知集合 $A=\{a\mid a=2^x3^y5^z,\ x,y,z\in\N\}$，$B=\{b\mid b\in A,\ 1\leqslant b\leqslant 10\}$，试求集合 $B$ 中含有元素 $1$ 和 $10$ 的真子集的个数为 \blankline[4.4em].

% 注：原卷第 8 题只印「【解析】64」，无过程。按 $B=\{1,2,3,4,5,6,8,9,10\}$（9 个元素）计算，
%     含 $1$ 与 $10$ 的真子集应为 $2^7-1=127$，与 $64$（$=2^6$）不符；原卷口径无法确证，
%     本稿如原卷照录 $64$，不代为改写。
\ans{$64$}

\item 已知 $A=\left\{x\mid \dfrac{x+a}{x^2-4}=1,\ x\in\R\right\}$ 仅有两个子集，则实数 $a$ 取值的集合为 \blankline[4.4em].

\ans{$\left\{\pm 2,-\dfrac{17}{4}\right\}$}

\ifanswer
\solbegin
\step{因为 $A=\{x\mid \dfrac{x+a}{x^2-4}=1,\ x\in\R\}$ 仅有两个子集，所以 $A$ 只有一个元素，$\dfrac{x+a}{x^2-4}=1$，}
\step{①$a\neq\pm 2$，得 $x^2-x-a-4=0$，所以此方程有两个相等的实数根，}
\step{得 $\Delta=1-4(-a-4)=0$，解得 $a=-\dfrac{17}{4}$. 经过验证满足条件.}
\step{②$a=-2$ 时，$\dfrac{x+a}{x^2-4}=1$，化为 $1=x+2$，解得 $x=-1$ 满足条件.}
% 注：原卷此处 ③ 印作「③$a=-2$ 时」（与 ② 同为 $a=-2$），但该行紧随的「化为 $1=x-2$，解得 $x=3$」
%     只在 $a=2$ 时成立（$a=-2$ 时 $(x-2)/(x^2-4)=1/(x+2)=1$ 解得 $x=-1$，即 ② 的情形），
%     且本段结论「$a$ 取值的集合为 $\{\pm 2,-\frac{17}{4}\}$」含 $+2$，显系笔误，按文意订正为 $a=2$.
\step{③$a=2$ 时，$\dfrac{x+a}{x^2-4}=1$，化为 $1=x-2$，解得 $x=3$ 满足条件.}
\step{综上，实数 $a$ 取值的集合为 $\left\{\pm 2,-\dfrac{17}{4}\right\}$.}
\solend

\fi

\item 设集合 $I=\{1,2,3,4,5\}$，若非空集合 $A$ 同时满足①$A\subseteq I$，②$|A|\leqslant \min(A)$（其中 $|A|$ 表示 $A$ 中元素的个数，$\min(A)$ 表示集合 $A$ 中最小元素），称集合 $A$ 为 $I$ 的一个好子集，$I$ 的所有好子集的个数为 \blankline[2.4em].

\ans{$12$}

\ifanswer
\solbegin
\step{①当 $|A|=1$，即集合 $A$ 中元素的个数为 $1$ 时，}
\step{$A$ 的可能情况为 $\{1\}$，$\{2\}$，$\{3\}$，$\{4\}$，$\{5\}$，}
\step{②当 $|A|=2$，即集合 $A$ 中元素的个数为 $2$ 时，}
\step{$A$ 的可能情况为 $\{2,3\}$，$\{2,4\}$，$\{2,5\}$，$\{3,4\}$，$\{3,5\}$，$\{4,5\}$，}
\step{③当 $|A|=3$，即集合 $A$ 中元素的个数为 $3$ 时，$A$ 的可能情况为 $\{3,4,5\}$，}
\step{所以 $I$ 的所有好子集的个数为 $12$.}
\solend

\fi

\end{enumerate}

\bigsec{二、选择题}

\begin{enumerate}
\setcounter{enumi}{10}

% 注：原卷第 11 题的 Venn 图在公众号原图里缺失（题干与选项之间只有正常行距，无任何图形像素），
%     图缺失即无从判断阴影区域，故本稿照录题干与四个选项，不补图、不代拟答案。
\item 下图中的阴影部分的正确表示是（\quad）

A. $(A\cup B)\cap(A\cup C)$\qquad B. $(A\cup B)\cap C$

C. $(A\cup C)\cap(B\cup C)$\qquad D. $(A\cup B)\cap(B\cup C)$

\ans{原卷此处应有三集合 Venn 图，但公众号原图缺失，无从判断阴影区域，故本稿不给出答案}

\item 有下列三个命题：①“$x+y\neq 4$”是“$x\neq 1$ 且 $y\neq 3$”的必要非充分条件；②$xy<0$ 是 $|x-y|=|x|+|y|$ 的充要条件；③已知 $m$、$n\in\Z$，则 $m^2+n^2<5$ 是 $m+n\leqslant 2$ 的充分不必要条件；其中的真命题有（\quad）

A. $0$ 个\qquad B. $1$ 个\qquad C. $2$ 个\qquad D. $3$ 个

\ans{$B$}

\ifanswer
\solbegin
\step{①“$x+y\neq 4$”，推不出“$x\neq 1$ 且 $y\neq 3$”，反之不成立，例如 $x=5$，$y=-1$，则 $x+y=4$，}
\step{因此“$x+y\neq 4$”是“$x\neq 1$ 且 $y\neq 3$”的既不充分也不必要条件，因此①是假命题，}
% 注：原卷 ② 一行印作「②$xy\leqslant 0$ 是 $|x-y|=|x|+|y|$ 的充要条件，因此②是假命题」。
%     但 $xy\leqslant 0$ 与 $|x-y|=|x|+|y|$ 确实互为充要条件，该命题应为**真**，与同行的
%     「因此②是假命题」以及本段末的「故选 B」（须恰有 1 个真命题，即只能 ③ 为真）都冲突；
%     而题干 ② 印的是 $xy<0$，$xy<0$ 只是充分不必要条件，恰为假命题。故本稿按题干与文意
%     把解析该行订正为「$xy<0$ 是 $|x-y|=|x|+|y|$ 的充分不必要条件」。
\step{②$xy<0$ 是 $|x-y|=|x|+|y|$ 的充分不必要条件，因此②是假命题；}
\step{③若 $m$、$n\in\Z$，$m^2+n^2<5$，则 $(m,n)$ 为 $(1,0)$，$(2,0)$，$(1,1)$，$(-1,0)$，$(-2,0)$，$(-1,1)$，$(0,1)$，$(0,2)$，$(0,-1)$，$(0,-2)$，$(-1,-1)$，$(1,-1)$，共 $9$ 组，}
\step{都满足 $m+n\leqslant 2$，反之不成立，例如：$m=3$，$n=-1$，}
\step{所以 $m$、$n\in\Z$，则 $m^2+n^2<5$ 是 $m+n\leqslant 2$ 的充分不必要条件，因此③是真命题.}
\step{故选 $B$.}
\solend

\fi

\item 若 $a$、$b$、$c$ 是 $\triangle ABC$ 的三条边，则“$a^2+b^2+c^2=ab+bc+ca$”是“$\triangle ABC$ 是等腰三角形”的（\quad）

A. 充分不必要条件\qquad B. 必要不充分条件

C. 充要条件\qquad D. 既不充分也不必要条件

\ans{$A$}

\ifanswer
\solbegin
\step{若“$\triangle ABC$ 是等腰三角形”，则当 $a=b\neq c$，}
\step{则 $a^2+b^2+c^2=ab+bc+ca$ 不一定成立，}
\step{若 $a^2+b^2+c^2=ab+bc+ca$，则 $2a^2+2b^2+2c^2=2ab+2bc+2ca$，}
\step{即 $(a-b)^2+(b-c)^2+(c-a)^2=0$，即 $a-b=0$，$b-c=0$，$c-a=0$，}
\step{则 $a=b=c$，则“$\triangle ABC$ 是等腰三角形”成立，}
\step{即“$a^2+b^2+c^2=ab+bc+ca$”是“$\triangle ABC$ 是等腰三角形”充分不必要条件，}
\step{故选 $A$.}
\solend

\fi

\item 设集合 $P_1=\{x\mid x^2+ax+1>0\}$，$P_2=\{x\mid x^2+ax+2>0\}(a\in\R)$，下列说法正确的是（\quad）

A. 对任意 $a$，$P_1$ 是 $P_2$ 的真子集\qquad B. 对任意 $a$，$P_1$ 不是 $P_2$ 的子集

C. 存在 $a$，使得 $P_1$ 不是 $P_2$ 的子集\qquad D. 存在 $a$，使得 $P_2$ 是 $P_1$ 的真子集

\ans{$A$}

\ifanswer
\solbegin
\step{由 $x^2+ax+1>0$ 得 $x^2+ax+2=x^2+ax+1+1>1>0$ 一定成立，}
\step{但当 $x^2+ax+2>0$ 时，$x^2+ax+1>0$ 不一定成立，}
\step{故对于任意的 $a$ 都有 $P_1\bsub P_2$. 故选 $A$.}
\solend

\fi

\end{enumerate}

\bigsec{三、解答题}

\begin{enumerate}
\setcounter{enumi}{14}

\item 求证：关于 $x$ 的方程 $ax^2+bx+c=0$ 有一根为 $1$ 的充要条件是 $a+b+c=0$.

\ans{略}

\ansspace[8]

\item 已知命题“关于 $x$ 的方程 $x^2+mx+2m+5=0$ 有两个不相等的实数根”是假命题.

\begin{enumerate}[label=(\arabic*)]
\item 求实数 $m$ 的取值集合 $A$；
\item 设集合 $B=\{x\mid 1-2a\leqslant x\leqslant a-1\}$，若 $x\in A$ 是 $x\in B$ 的充分不必要条件，求实数 $a$ 的取值范围.
\end{enumerate}

\ans{（1）$A=\{m\mid -2\leqslant m\leqslant 10\}$；（2）$[11,+\infty)$}

\ifanswer
\solbegin
% 注：原卷 (1) 的解析印作「若命题……是真命题，所以 $\Delta=m^2-4(2m+5)>0$，解得 $m>10$ 或 $m<-2$」，
%     与题干「……是假命题」以及紧随的结论「故 $A=\{m\mid -2\leqslant m\leqslant 10\}$」都矛盾
%     （按 $\Delta>0$ 应得 $m>10$ 或 $m<-2$），显系笔误；本稿按文意订正为「是假命题」「$\Delta\leqslant 0$」
%     「$-2\leqslant m\leqslant 10$」，结论不变。
\stepm{(1)}{若命题“关于 $x$ 的方程 $x^2+mx+2m+5=0$ 有两个不相等的实数根”是假命题，所以 $\Delta=m^2-4(2m+5)\leqslant 0$，解得 $-2\leqslant m\leqslant 10$，}
\step{故 $A=\{m\mid -2\leqslant m\leqslant 10\}$.}
\stepm{(2)}{因为 $A=\{m\mid -2\leqslant m\leqslant 10\}$，$x\in A$ 是 $x\in B$ 的充分不必要条件，所以 $A\subset B$.}
\step{即 $\begin{cases}1-2a\leqslant -2\\ a-1\geqslant 10\end{cases}$ 且等号不同时成立，解得 $a\geqslant 11$，}
\step{所以实数 $a$ 的取值范围为 $[11,+\infty)$.}
\solend

\fi

\ansspace[6]

\item 设非空集合 $S=\{x\mid m\leqslant x\leqslant n\}$ 满足：当 $x\in S$ 时，有 $x^2\in S$.

\begin{enumerate}[label=(\arabic*)]
\item 当 $m=1$ 时，求 $S$；
\item 若 $m=-\dfrac12$，求实数 $n$ 的取值范围.
\end{enumerate}

\ans{（1）$S=\{1\}$；（2）$\left[\dfrac14,1\right]$}

\ifanswer
\solbegin
\stepm{(1)}{若 $m=1$ 时，所以 $S=\{x\mid 1\leqslant x\leqslant n\}$，所以 $\forall x$ 有 $1\leqslant x^2\leqslant n^2$，}
\step{所以 $n\leqslant n^2$，所以 $n\leqslant 1$，所以 $S=\{1\}$；}
\stepm{(2)}{若 $m=-\dfrac12$，}
\step{①$-\dfrac12\leqslant n\leqslant 0$ 时，$\forall x\in S$，有 $n^2\leqslant x^2\leqslant \dfrac14$，所以 $x^2\notin S$，舍去，}
\step{②$0<n\leqslant \dfrac12$ 时，$\forall x\in S$，有 $0\leqslant x^2\leqslant \dfrac14$，所以 $\dfrac14\leqslant n\leqslant \dfrac12$，}
\step{③$n>\dfrac12$ 时，$\forall x\in S$，有 $0\leqslant x^2\leqslant n^2$，所以 $n^2\leqslant n$，故 $\dfrac12<n\leqslant 1$.}
\solend

\fi

\ansspace[6]

% 压轴题：其后已无内容，故作答区全用可丢弃胶水（\ansspace*），并在题目前先量一下版面。
\ansneed{7cm}

\item 已知集合 $A=\{x\mid x=m^2-n^2,\ m$、$n\in\Z\}$.

\begin{enumerate}[label=(\arabic*)]
\item 判断 $8$，$9$，$10$ 是否属于集合 $A$；
\item 已知集合 $B=\{x\mid x=2k+1,k\in\Z\}$，证明：“$x\in A$”是“$x\in B$”的必要非充分条件；
\item 求所有满足集合 $A$ 的偶数，并说明理由.
\end{enumerate}

\ans{（1）$8\in A$，$9\in A$，$10\notin A$；（2）略；（3）所有满足集合 $A$ 的偶数为 $4k(k\in\Z)$}

\ifanswer
\solbegin
\stepm{(1)}{由于 $8=3^2-1^2$，满足集合 $A$ 中元素特征 $x=m^2-n^2$，$m$、$n\in\Z$，所以 $8\in A$，}
\step{由于 $9=5^2-4^2$，满足集合 $A$ 中元素特征 $x=m^2-n^2$，$m$、$n\in\Z$，所以 $9\in A$，}
\step{假设 $10=m^2-n^2$，$m$、$n\in\Z$，}
\step{则 $(|m|+|n|)(|m|-|n|)=10$，且 $|m|+|n|>|m|-|n|>0$，}
% 注：原卷此行印作「由于 $10=1\times10\times2\times5$」（三个乘号），右端为 100 而非 10；
%     紧随的两组方程组正是按 $1\times10$ 与 $2\times5$ 两个分解写出的，显系漏排等号，
%     本稿按文意订正为「$10=1\times10=2\times5$」（详见文件头「原卷笔误」④）。
\step{由于 $10=1\times 10=2\times 5$，所以 $\begin{cases}|m|+|n|=10\\ |m|-|n|=1\end{cases}$ 或 $\begin{cases}|m|+|n|=5\\ |m|-|n|=2\end{cases}$，}
\step{显然均无整数解，所以 $10\notin A$；}
\stepm{(2)}{集合 $B=\{x\mid x=2k+1,k\in\Z\}$，则恒有 $2k+1=(k+1)^2-k^2$，}
\step{即满足集合 $A$ 中元素特征 $x=m^2-n^2$，$m$、$n\in\Z$，所以 $2k+1\in A$，}
\step{即一切奇数都属于 $A$；}
\step{反之满足这个集合中的元素不一定全是奇数，如 $8\in A$，}
\step{所以 $x\in A$ 是 $x\in B$ 的必要非充分条件，}
\stepm{(3)}{集合 $A=\{x\mid x=m^2-n^2\}(m,n\in\Z)$，而 $m^2-n^2=(m+n)(m-n)$，}
\step{①当 $m$ 和 $n$ 一奇一偶时，$m+n$ 和 $m-n$ 均为奇数，}
\step{所以 $(m+n)(m-n)$ 为奇数，不满足题意；}
\step{②当 $m$ 和 $n$ 同为奇数和偶数时，$m-n$，$m+n$ 均为偶数，}
\step{所以 $(m+n)(m-n)$ 为 $4$ 的倍数，}
% 注：原卷此处印作「反之当 $B=\{x\mid x=2k+1,k\in\Z\}$」（与第（2）问定义的 $B$ 一字不差），
%     但紧随的推导「不妨令 $m+n=2k$，$m-n=2$」所得正是 $x=(m+n)(m-n)=4k$，
%     且本段结论为「所有满足集合 $A$ 的偶数为 $4k(k\in\Z)$」；此处 $B$ 疑为 $\{x\mid x=4k,k\in\Z\}$
%     之误（或漏改自第（2）问），但无法确证，本稿照抄原卷写法。
\step{反之当 $B=\{x\mid x=2k+1,k\in\Z\}$，则不妨令 $\begin{cases}m+n=2k\\ m-n=2\end{cases}$，}
\step{可解得 $\begin{cases}m=k+1\\ n=k-1\end{cases}$，满足集合 $A$ 中元素特征 $x=m^2-n^2$，$m$、$n\in\Z$，}
\step{所以满足集合 $A$ 的偶数为 $4k(k\in\Z)$；}
\step{综上所述，所有满足集合 $A$ 的偶数为 $4k(k\in\Z)$.}
\solend

\fi

\ansspace*[10]

\end{enumerate}

\end{document}
"
%
% 原始材料是微信公众号图文（公众号「破晓」，原文链接 https://mp.weixin.qq.com/s/GcU0DDwonIGn8iVGwVRmpg），
% 正文为 7 张 PNG 页；原文本身就是一份**讲评版**——题目下方直接印出【答案】或【解析】。
% 试题文字、答案与解析均照原卷录入，试卷版把所有行间答案改排 \ans{}、解析收进 \ifanswer…\fi 块。
%
% 与原卷的排版差异（均已在 README「说明」中交代）：
%   ① 原文自**第 3 题**起（第 1、2 题未在该文中发布），故本稿题号亦自 3 起，与高一03（同校作业 1）一致；
%   ② 原卷本就印有「一、填空题」「二、选择题」「三、解答题」三节标题（分见原图 00/02/03.png），
%      本稿照原卷分节，只把标题改排成全稿统一的「大节标题」样式；
%   ③ 原卷把答案直接印在题目下方（选择题第 12–14 题印在【解析】里），本稿统一改为试卷版隐去、
%      答案版以【答案】给出，使两版彻底分离；选择题的括注一律改排「（\quad）」；
%   ④ 原卷的实数集、自然数集符号印作普通斜体 $R$、$N$、$N^*$（第 9、10、12、14、18 题等），
%      本稿按全稿惯例统一写作 $\R$、$\N$、$\N^*$；原卷中文括注用**半角**括号（第 11–14 题的
%      「(   )」），本稿按全稿惯例排全角；
%   ⑤ 第 11 题原卷题干下方本应有一幅三集合 Venn 图，但公众号原图该处为空白（题干与选项之间
%      只有正常行距，逐像素核对过），图确实缺失。本稿照录题干、不补图、不代拟答案，详见该题下的注。
%
% 原卷笔误四处（已逐处放大到能看清字形后核对，本稿按文意订正，并在 README 中写明原委）：
%   ① 第 9 题【解析】③ 原卷印作「③$a=-2$ 时」（与 ② 同为 $a=-2$），但该行紧随的
%      「化为 $1=x-2$，解得 $x=3$」只在 $a=2$ 时成立，且本段结论「$a$ 取值的集合为
%      $\{\pm 2,-\frac{17}{4}\}$」含 $+2$，显系笔误，本稿订正为「③$a=2$ 时」；
%   ② 第 16 题【解析】(1) 原卷印作「若命题……是真命题，所以 $\Delta=m^2-4(2m+5)>0$，
%      解得 $m>10$ 或 $m<-2$」，与题干「是假命题」及紧随的结论「故 $A=\{m\mid -2\leqslant m\leqslant 10\}$」
%      都矛盾（按 $\Delta>0$ 应得 $m>10$ 或 $m<-2$），按文意订正为「是假命题」「$\Delta\leqslant 0$」
%      「$-2\leqslant m\leqslant 10$」（结论不变）。
%   ③ 第 12 题【解析】② 原卷印作「$xy\leqslant 0$ 是 $|x-y|=|x|+|y|$ 的充要条件，因此②是假命题」，
%      但 $xy\leqslant 0$ 与 $|x-y|=|x|+|y|$ 确为充要条件、该命题应为真，与同行的「因此②是假命题」
%      及段末「故选 B」（须恰有 1 个真命题）冲突；而题干 ② 印的正是 $xy<0$（充分不必要条件，恰为假），
%      故本稿按题干与文意把该行订正为「$xy<0$ 是 $|x-y|=|x|+|y|$ 的充分不必要条件」；
%   ④ 第 18 题【解析】(1) 原卷印作「由于 $10=1\times10\times2\times5$」（三个乘号），但该式左端为 10
%      而右端为 100；紧随的两组方程组恰是按 $1\times10$ 与 $2\times5$ 这两个分解写出的，显系漏排等号，
%      本稿订正为「$10=1\times10=2\times5$」。
% 另：第 16 题(2)【解析】的「所以 $A\subset B$」照原卷录入（原卷为**真包含**，与第 10 题①的
%     $A\subseteq I$ 逐处放大比对后确认不同）；2026-09-15 回原图复核时曾误作 $\subseteq$，已订正。
% 原卷存疑两处（无法确证，本稿照抄并加注）：
%   ① 第 8 题只印「【解析】64」；② 第 18(3)【解析】中「反之当 $B=\{x\mid x=2k+1,k\in\Z\}$」。

\documentclass[a4paper,11pt]{article}
\usepackage{common}

\begin{document}

\examtitlest[3]{2026--2027 年格致中学高一上周末作业 2}{集合与常用逻辑用语}

\bigsec{一、填空题}

\begin{enumerate}
\setcounter{enumi}{2}

\item 若非空集合 $A=\{x\mid -2m+6<x<m-2\}$，$B=\{x\mid -m<x<m\}$，且 $A\subset B$，则实数 $m$ 的取值范围是 \blankline[4.4em].

\ans{$\left(\dfrac83,6\right]$}

\item 已知条件 $p:2k-1\leqslant x\leqslant -3k$，条件 $q:-1<x<3$，$p$ 是 $q$ 的必要条件，则实数 $k$ 的取值范围为 \blankline[4.4em].

\ans{$(-\infty,-1]$}

\ifanswer
\solbegin
\step{若 $p$ 是 $q$ 的必要条件，则 $q\Rightarrow p$，即 $\begin{cases}-3k\geqslant 3\\ 2k-1\leqslant -1\end{cases}$，得 $\begin{cases}k\leqslant -1\\ k\leqslant 0\end{cases}$，得 $k\leqslant -1$，}
\step{即实数 $k$ 的取值范围是 $(-\infty,-1]$.}
\solend

\fi

\item 集合 $A=\{0,1,a\}$，$B=(0,2)$，若 $A\cap B=\{1\}$，则实数 $a$ 的取值范围是 \blankline[4.4em].

\ans{$(-\infty,0]\cup[2,+\infty)$}

\item 已知 $\varnothing\subset M\subseteq\N$，且当 $x\in M$ 时，$6-x\in M$，若集合 $M$ 只有三个元素，这样的集合 $M$ 为 \blankline[4.4em].

% 注：原卷第 6 题此处印作「$\varnothing\subset M\subseteq N$」，其中「$N$」为普通斜体 $N$（无空心双竖线）。
%     由答案 $\{0,3,6\}$、$\{1,3,5\}$、$\{2,3,4\}$ 反推，$6-x\in M$ 须把元素限制在 $[0,6]$ 内，
%     故该符号即自然数集 $\N$，本稿按全稿惯例写作 $\N$。
\ans{$\{0,3,6\}$ 或 $\{1,3,5\}$ 或 $\{2,3,4\}$}

\item 任意两个正整数 $x$、$y$，定义某种运算 $\otimes$：$x\otimes y=\begin{cases}x+y & (x\text{ 与 }y\text{ 奇偶相同})\\ x\times y & (x\text{ 与 }y\text{ 奇偶不同})\end{cases}$，则集合 $M=\{(x,y)\mid x\otimes y=6,\ x,y\in\N^*\}$ 中元素的个数是 \blankline[2.4em].

\ans{$9$}

\ifanswer
\solbegin
\step{①当 $x$ 与 $y$ 都为奇数时，有 $1+5=6$，$3+3=6$，}
\step{得 $(1,5)$、$(5,1)$、$(3,3)$，$3$ 个点符合题意，}
\step{②当 $x$ 与 $y$ 都为偶数时，有 $2+4=6$，}
\step{得 $(2,4)$、$(4,2)$，$2$ 个点符合题意，}
\step{③当 $x$ 与 $y$ 一奇一偶时，$1\times 6=6$，$2\times 3=6$，}
\step{得 $(1,6)$、$(6,1)$、$(2,3)$、$(3,2)$，$4$ 个点符合题意，}
\step{所以共有 $9$ 个点符合题意.}
\solend

\fi

\item 已知集合 $A=\{a\mid a=2^x3^y5^z,\ x,y,z\in\N\}$，$B=\{b\mid b\in A,\ 1\leqslant b\leqslant 10\}$，试求集合 $B$ 中含有元素 $1$ 和 $10$ 的真子集的个数为 \blankline[4.4em].

% 注：原卷第 8 题只印「【解析】64」，无过程。按 $B=\{1,2,3,4,5,6,8,9,10\}$（9 个元素）计算，
%     含 $1$ 与 $10$ 的真子集应为 $2^7-1=127$，与 $64$（$=2^6$）不符；原卷口径无法确证，
%     本稿如原卷照录 $64$，不代为改写。
\ans{$64$}

\item 已知 $A=\left\{x\mid \dfrac{x+a}{x^2-4}=1,\ x\in\R\right\}$ 仅有两个子集，则实数 $a$ 取值的集合为 \blankline[4.4em].

\ans{$\left\{\pm 2,-\dfrac{17}{4}\right\}$}

\ifanswer
\solbegin
\step{因为 $A=\{x\mid \dfrac{x+a}{x^2-4}=1,\ x\in\R\}$ 仅有两个子集，所以 $A$ 只有一个元素，$\dfrac{x+a}{x^2-4}=1$，}
\step{①$a\neq\pm 2$，得 $x^2-x-a-4=0$，所以此方程有两个相等的实数根，}
\step{得 $\Delta=1-4(-a-4)=0$，解得 $a=-\dfrac{17}{4}$. 经过验证满足条件.}
\step{②$a=-2$ 时，$\dfrac{x+a}{x^2-4}=1$，化为 $1=x+2$，解得 $x=-1$ 满足条件.}
% 注：原卷此处 ③ 印作「③$a=-2$ 时」（与 ② 同为 $a=-2$），但该行紧随的「化为 $1=x-2$，解得 $x=3$」
%     只在 $a=2$ 时成立（$a=-2$ 时 $(x-2)/(x^2-4)=1/(x+2)=1$ 解得 $x=-1$，即 ② 的情形），
%     且本段结论「$a$ 取值的集合为 $\{\pm 2,-\frac{17}{4}\}$」含 $+2$，显系笔误，按文意订正为 $a=2$.
\step{③$a=2$ 时，$\dfrac{x+a}{x^2-4}=1$，化为 $1=x-2$，解得 $x=3$ 满足条件.}
\step{综上，实数 $a$ 取值的集合为 $\left\{\pm 2,-\dfrac{17}{4}\right\}$.}
\solend

\fi

\item 设集合 $I=\{1,2,3,4,5\}$，若非空集合 $A$ 同时满足①$A\subseteq I$，②$|A|\leqslant \min(A)$（其中 $|A|$ 表示 $A$ 中元素的个数，$\min(A)$ 表示集合 $A$ 中最小元素），称集合 $A$ 为 $I$ 的一个好子集，$I$ 的所有好子集的个数为 \blankline[2.4em].

\ans{$12$}

\ifanswer
\solbegin
\step{①当 $|A|=1$，即集合 $A$ 中元素的个数为 $1$ 时，}
\step{$A$ 的可能情况为 $\{1\}$，$\{2\}$，$\{3\}$，$\{4\}$，$\{5\}$，}
\step{②当 $|A|=2$，即集合 $A$ 中元素的个数为 $2$ 时，}
\step{$A$ 的可能情况为 $\{2,3\}$，$\{2,4\}$，$\{2,5\}$，$\{3,4\}$，$\{3,5\}$，$\{4,5\}$，}
\step{③当 $|A|=3$，即集合 $A$ 中元素的个数为 $3$ 时，$A$ 的可能情况为 $\{3,4,5\}$，}
\step{所以 $I$ 的所有好子集的个数为 $12$.}
\solend

\fi

\end{enumerate}

\bigsec{二、选择题}

\begin{enumerate}
\setcounter{enumi}{10}

% 注：原卷第 11 题的 Venn 图在公众号原图里缺失（题干与选项之间只有正常行距，无任何图形像素），
%     图缺失即无从判断阴影区域，故本稿照录题干与四个选项，不补图、不代拟答案。
\item 下图中的阴影部分的正确表示是（\quad）

A. $(A\cup B)\cap(A\cup C)$\qquad B. $(A\cup B)\cap C$

C. $(A\cup C)\cap(B\cup C)$\qquad D. $(A\cup B)\cap(B\cup C)$

\ans{原卷此处应有三集合 Venn 图，但公众号原图缺失，无从判断阴影区域，故本稿不给出答案}

\item 有下列三个命题：①“$x+y\neq 4$”是“$x\neq 1$ 且 $y\neq 3$”的必要非充分条件；②$xy<0$ 是 $|x-y|=|x|+|y|$ 的充要条件；③已知 $m$、$n\in\Z$，则 $m^2+n^2<5$ 是 $m+n\leqslant 2$ 的充分不必要条件；其中的真命题有（\quad）

A. $0$ 个\qquad B. $1$ 个\qquad C. $2$ 个\qquad D. $3$ 个

\ans{$B$}

\ifanswer
\solbegin
\step{①“$x+y\neq 4$”，推不出“$x\neq 1$ 且 $y\neq 3$”，反之不成立，例如 $x=5$，$y=-1$，则 $x+y=4$，}
\step{因此“$x+y\neq 4$”是“$x\neq 1$ 且 $y\neq 3$”的既不充分也不必要条件，因此①是假命题，}
% 注：原卷 ② 一行印作「②$xy\leqslant 0$ 是 $|x-y|=|x|+|y|$ 的充要条件，因此②是假命题」。
%     但 $xy\leqslant 0$ 与 $|x-y|=|x|+|y|$ 确实互为充要条件，该命题应为**真**，与同行的
%     「因此②是假命题」以及本段末的「故选 B」（须恰有 1 个真命题，即只能 ③ 为真）都冲突；
%     而题干 ② 印的是 $xy<0$，$xy<0$ 只是充分不必要条件，恰为假命题。故本稿按题干与文意
%     把解析该行订正为「$xy<0$ 是 $|x-y|=|x|+|y|$ 的充分不必要条件」。
\step{②$xy<0$ 是 $|x-y|=|x|+|y|$ 的充分不必要条件，因此②是假命题；}
\step{③若 $m$、$n\in\Z$，$m^2+n^2<5$，则 $(m,n)$ 为 $(1,0)$，$(2,0)$，$(1,1)$，$(-1,0)$，$(-2,0)$，$(-1,1)$，$(0,1)$，$(0,2)$，$(0,-1)$，$(0,-2)$，$(-1,-1)$，$(1,-1)$，共 $9$ 组，}
\step{都满足 $m+n\leqslant 2$，反之不成立，例如：$m=3$，$n=-1$，}
\step{所以 $m$、$n\in\Z$，则 $m^2+n^2<5$ 是 $m+n\leqslant 2$ 的充分不必要条件，因此③是真命题.}
\step{故选 $B$.}
\solend

\fi

\item 若 $a$、$b$、$c$ 是 $\triangle ABC$ 的三条边，则“$a^2+b^2+c^2=ab+bc+ca$”是“$\triangle ABC$ 是等腰三角形”的（\quad）

A. 充分不必要条件\qquad B. 必要不充分条件

C. 充要条件\qquad D. 既不充分也不必要条件

\ans{$A$}

\ifanswer
\solbegin
\step{若“$\triangle ABC$ 是等腰三角形”，则当 $a=b\neq c$，}
\step{则 $a^2+b^2+c^2=ab+bc+ca$ 不一定成立，}
\step{若 $a^2+b^2+c^2=ab+bc+ca$，则 $2a^2+2b^2+2c^2=2ab+2bc+2ca$，}
\step{即 $(a-b)^2+(b-c)^2+(c-a)^2=0$，即 $a-b=0$，$b-c=0$，$c-a=0$，}
\step{则 $a=b=c$，则“$\triangle ABC$ 是等腰三角形”成立，}
\step{即“$a^2+b^2+c^2=ab+bc+ca$”是“$\triangle ABC$ 是等腰三角形”充分不必要条件，}
\step{故选 $A$.}
\solend

\fi

\item 设集合 $P_1=\{x\mid x^2+ax+1>0\}$，$P_2=\{x\mid x^2+ax+2>0\}(a\in\R)$，下列说法正确的是（\quad）

A. 对任意 $a$，$P_1$ 是 $P_2$ 的真子集\qquad B. 对任意 $a$，$P_1$ 不是 $P_2$ 的子集

C. 存在 $a$，使得 $P_1$ 不是 $P_2$ 的子集\qquad D. 存在 $a$，使得 $P_2$ 是 $P_1$ 的真子集

\ans{$A$}

\ifanswer
\solbegin
\step{由 $x^2+ax+1>0$ 得 $x^2+ax+2=x^2+ax+1+1>1>0$ 一定成立，}
\step{但当 $x^2+ax+2>0$ 时，$x^2+ax+1>0$ 不一定成立，}
\step{故对于任意的 $a$ 都有 $P_1\bsub P_2$. 故选 $A$.}
\solend

\fi

\end{enumerate}

\bigsec{三、解答题}

\begin{enumerate}
\setcounter{enumi}{14}

\item 求证：关于 $x$ 的方程 $ax^2+bx+c=0$ 有一根为 $1$ 的充要条件是 $a+b+c=0$.

\ans{略}

\ansspace[8]

\item 已知命题“关于 $x$ 的方程 $x^2+mx+2m+5=0$ 有两个不相等的实数根”是假命题.

\begin{enumerate}[label=(\arabic*)]
\item 求实数 $m$ 的取值集合 $A$；
\item 设集合 $B=\{x\mid 1-2a\leqslant x\leqslant a-1\}$，若 $x\in A$ 是 $x\in B$ 的充分不必要条件，求实数 $a$ 的取值范围.
\end{enumerate}

\ans{（1）$A=\{m\mid -2\leqslant m\leqslant 10\}$；（2）$[11,+\infty)$}

\ifanswer
\solbegin
% 注：原卷 (1) 的解析印作「若命题……是真命题，所以 $\Delta=m^2-4(2m+5)>0$，解得 $m>10$ 或 $m<-2$」，
%     与题干「……是假命题」以及紧随的结论「故 $A=\{m\mid -2\leqslant m\leqslant 10\}$」都矛盾
%     （按 $\Delta>0$ 应得 $m>10$ 或 $m<-2$），显系笔误；本稿按文意订正为「是假命题」「$\Delta\leqslant 0$」
%     「$-2\leqslant m\leqslant 10$」，结论不变。
\stepm{(1)}{若命题“关于 $x$ 的方程 $x^2+mx+2m+5=0$ 有两个不相等的实数根”是假命题，所以 $\Delta=m^2-4(2m+5)\leqslant 0$，解得 $-2\leqslant m\leqslant 10$，}
\step{故 $A=\{m\mid -2\leqslant m\leqslant 10\}$.}
\stepm{(2)}{因为 $A=\{m\mid -2\leqslant m\leqslant 10\}$，$x\in A$ 是 $x\in B$ 的充分不必要条件，所以 $A\subset B$.}
\step{即 $\begin{cases}1-2a\leqslant -2\\ a-1\geqslant 10\end{cases}$ 且等号不同时成立，解得 $a\geqslant 11$，}
\step{所以实数 $a$ 的取值范围为 $[11,+\infty)$.}
\solend

\fi

\ansspace[6]

\item 设非空集合 $S=\{x\mid m\leqslant x\leqslant n\}$ 满足：当 $x\in S$ 时，有 $x^2\in S$.

\begin{enumerate}[label=(\arabic*)]
\item 当 $m=1$ 时，求 $S$；
\item 若 $m=-\dfrac12$，求实数 $n$ 的取值范围.
\end{enumerate}

\ans{（1）$S=\{1\}$；（2）$\left[\dfrac14,1\right]$}

\ifanswer
\solbegin
\stepm{(1)}{若 $m=1$ 时，所以 $S=\{x\mid 1\leqslant x\leqslant n\}$，所以 $\forall x$ 有 $1\leqslant x^2\leqslant n^2$，}
\step{所以 $n\leqslant n^2$，所以 $n\leqslant 1$，所以 $S=\{1\}$；}
\stepm{(2)}{若 $m=-\dfrac12$，}
\step{①$-\dfrac12\leqslant n\leqslant 0$ 时，$\forall x\in S$，有 $n^2\leqslant x^2\leqslant \dfrac14$，所以 $x^2\notin S$，舍去，}
\step{②$0<n\leqslant \dfrac12$ 时，$\forall x\in S$，有 $0\leqslant x^2\leqslant \dfrac14$，所以 $\dfrac14\leqslant n\leqslant \dfrac12$，}
\step{③$n>\dfrac12$ 时，$\forall x\in S$，有 $0\leqslant x^2\leqslant n^2$，所以 $n^2\leqslant n$，故 $\dfrac12<n\leqslant 1$.}
\solend

\fi

\ansspace[6]

% 压轴题：其后已无内容，故作答区全用可丢弃胶水（\ansspace*），并在题目前先量一下版面。
\ansneed{7cm}

\item 已知集合 $A=\{x\mid x=m^2-n^2,\ m$、$n\in\Z\}$.

\begin{enumerate}[label=(\arabic*)]
\item 判断 $8$，$9$，$10$ 是否属于集合 $A$；
\item 已知集合 $B=\{x\mid x=2k+1,k\in\Z\}$，证明：“$x\in A$”是“$x\in B$”的必要非充分条件；
\item 求所有满足集合 $A$ 的偶数，并说明理由.
\end{enumerate}

\ans{（1）$8\in A$，$9\in A$，$10\notin A$；（2）略；（3）所有满足集合 $A$ 的偶数为 $4k(k\in\Z)$}

\ifanswer
\solbegin
\stepm{(1)}{由于 $8=3^2-1^2$，满足集合 $A$ 中元素特征 $x=m^2-n^2$，$m$、$n\in\Z$，所以 $8\in A$，}
\step{由于 $9=5^2-4^2$，满足集合 $A$ 中元素特征 $x=m^2-n^2$，$m$、$n\in\Z$，所以 $9\in A$，}
\step{假设 $10=m^2-n^2$，$m$、$n\in\Z$，}
\step{则 $(|m|+|n|)(|m|-|n|)=10$，且 $|m|+|n|>|m|-|n|>0$，}
% 注：原卷此行印作「由于 $10=1\times10\times2\times5$」（三个乘号），右端为 100 而非 10；
%     紧随的两组方程组正是按 $1\times10$ 与 $2\times5$ 两个分解写出的，显系漏排等号，
%     本稿按文意订正为「$10=1\times10=2\times5$」（详见文件头「原卷笔误」④）。
\step{由于 $10=1\times 10=2\times 5$，所以 $\begin{cases}|m|+|n|=10\\ |m|-|n|=1\end{cases}$ 或 $\begin{cases}|m|+|n|=5\\ |m|-|n|=2\end{cases}$，}
\step{显然均无整数解，所以 $10\notin A$；}
\stepm{(2)}{集合 $B=\{x\mid x=2k+1,k\in\Z\}$，则恒有 $2k+1=(k+1)^2-k^2$，}
\step{即满足集合 $A$ 中元素特征 $x=m^2-n^2$，$m$、$n\in\Z$，所以 $2k+1\in A$，}
\step{即一切奇数都属于 $A$；}
\step{反之满足这个集合中的元素不一定全是奇数，如 $8\in A$，}
\step{所以 $x\in A$ 是 $x\in B$ 的必要非充分条件，}
\stepm{(3)}{集合 $A=\{x\mid x=m^2-n^2\}(m,n\in\Z)$，而 $m^2-n^2=(m+n)(m-n)$，}
\step{①当 $m$ 和 $n$ 一奇一偶时，$m+n$ 和 $m-n$ 均为奇数，}
\step{所以 $(m+n)(m-n)$ 为奇数，不满足题意；}
\step{②当 $m$ 和 $n$ 同为奇数和偶数时，$m-n$，$m+n$ 均为偶数，}
\step{所以 $(m+n)(m-n)$ 为 $4$ 的倍数，}
% 注：原卷此处印作「反之当 $B=\{x\mid x=2k+1,k\in\Z\}$」（与第（2）问定义的 $B$ 一字不差），
%     但紧随的推导「不妨令 $m+n=2k$，$m-n=2$」所得正是 $x=(m+n)(m-n)=4k$，
%     且本段结论为「所有满足集合 $A$ 的偶数为 $4k(k\in\Z)$」；此处 $B$ 疑为 $\{x\mid x=4k,k\in\Z\}$
%     之误（或漏改自第（2）问），但无法确证，本稿照抄原卷写法。
\step{反之当 $B=\{x\mid x=2k+1,k\in\Z\}$，则不妨令 $\begin{cases}m+n=2k\\ m-n=2\end{cases}$，}
\step{可解得 $\begin{cases}m=k+1\\ n=k-1\end{cases}$，满足集合 $A$ 中元素特征 $x=m^2-n^2$，$m$、$n\in\Z$，}
\step{所以满足集合 $A$ 的偶数为 $4k(k\in\Z)$；}
\step{综上所述，所有满足集合 $A$ 的偶数为 $4k(k\in\Z)$.}
\solend

\fi

\ansspace*[10]

\end{enumerate}

\end{document}
